% ===== 9 局限性、失败分析与改进方向 =====
\section{局限性、失败分析与改进方向}
\label{sec:limitations}

本文采用“预训练实例感知 + 显式物理尺度 + 低维筛分语义映射”的路线，优势是现场无需重新训练即可形成可解释输出，但这并不意味着单目视觉能够无条件替代机械筛分。系统的误差来自三个不同层次：有些是当前实现尚未充分校准，可以通过增加数据和工程修正降低；有些来自实例分割模型的领域失效，需要针对困难样本改进感知；还有一些属于单目二维观测本身缺失的信息，不能仅靠增加后处理公式消除。本章按“现象/边界--根因--影响--缓解”的结构主动说明这些限制，并据此确定后续工作的优先级。

\subsection{遮挡与不可见颗粒：单目观测的根本边界}
\label{sec:limit-occlusion}

\textbf{现象与边界。} 在自然堆积场景中，下层颗粒可能只露出局部轮廓，甚至完全被上层颗粒遮挡。实例分割只能对图像中实际可见的边界进行推断；当一个颗粒的大部分面积不可见时，即使模型正确识别出“这里存在一个实例”，由可见轮廓拟合得到的 $a_i$、$b_i$ 和面积也不再代表完整颗粒尺寸。完全被遮挡的颗粒更不存在任何单目视觉证据。

\textbf{根因。} 这不是 Cellpose-SAM 特有的问题，而是二维单视角观测缺少深度和背面信息的几何不可辨识性。任何只输入一张 RGB 图像的方法都无法从没有成像证据的位置唯一恢复真实颗粒数量和三维尺寸；端到端网络即使给出数值，也只能依赖训练集统计先验进行猜测。

\textbf{影响。} 局部遮挡通常导致颗粒轮廓和轴长低估，使个体 $d_i$ 偏小；多个颗粒边界被遮挡/阴影连接时还可能发生 merge，使单个伪实例尺寸偏大。前者推动级配偏细，后者在 $d^\gamma$ 质量权重下可能被放大并推动高分位数偏粗。因此，遮挡既可能产生相反方向的偏差，也会降低基于可见颗粒推断总体批次的代表性。

\textbf{缓解与改进。} 现场层面应优先采用稳定俯视、漫射照明，并在规则允许时减少严重二层堆叠；算法层面可在人工实例标注集上分别统计不同遮挡等级的 merge/split 与粒径误差，建立可用性阈值。若比赛场景必须保留高堆叠状态，真正的结构性改进是引入多视角、深度相机或三维重建，而不是继续调节二维筛分参数。

\subsection{空间尺度、透视与光学误差}
\label{sec:limit-scale}

\textbf{现象与边界。} 当前主链使用单一 $r$（mm/px）把整张图像从像素转换到毫米。如果相机光轴倾斜、镜头畸变明显、承载面不平或标尺与骨料不共面，同一像素长度在画面不同区域实际上对应不同物理长度，单一尺度只能给出近似。

\textbf{根因。} 式\eqref{eq:scale-conversion}隐含了近似正射投影和平面共面的成像假设。该假设在固定高位俯拍、视场相对较小时通常可以工程上逼近，但不是普适物理定律。当前 \texttt{infer\_resolution} 只从已经存在的 px/mm 双列反推尺度，也不会自动修复透视或镜头畸变。

\textbf{影响。} 尺度误差会近似按比例传递到 $a_i,b_i,d_{\mathrm{eq},i}$ 和 $d_i$，从而同时移动所有 $D_p$ 和筛级边界归属。与局部分割错误不同，尺度错误具有场景级系统性：如果 $r$ 整体偏大 5\%，大量颗粒可能共同向更粗筛级移动，因此它应在任何模型参数校准之前首先被控制。

\textbf{缓解与改进。} 当前部署应使用与测量面共面的参考物，并在画面多个位置交叉检查尺度。下一版本优先加入相机内参畸变校正和平面多点标记单应变换，将“一个场景一个 mm/px”升级为经几何校正后的统一测量平面。如果后续需要处理明显三维起伏堆体，则进一步采用深度或多视角，而不能用二维单应替代真实深度。

\subsection{二维视觉粒径与机械筛孔语义仍需物理校准}
\label{sec:limit-sieve-semantics}

\textbf{现象与边界。} 当前默认 $d=b$ 能够生成稳定的视觉粒径代理，但尚没有证据证明对不同来源、不同形貌骨料都与机械筛分通过行为一致。扁平颗粒、细长颗粒和近球形颗粒即使俯视短轴相同，其第三维厚度和在振筛中的可转动姿态也可能不同。

\textbf{根因。} 机械筛分是三维颗粒在动态筛孔中的通过问题，单目图像只是二维投影。式\eqref{eq:sieve-equivalent}用 $b$ 与 $d_{\mathrm{eq}}$ 的低维组合吸收平均系统差异，但无法从理论上恢复每个颗粒的真实筛孔通过方向。

\textbf{影响。} 如果未校准参数存在系统偏差，实例分割即使非常准确，最终 $D_{10}/D_{50}/D_{90}$ 和逐筛级质量占比仍可能整体偏离机械筛分。该误差直接作用于比赛主要评分对象，因此其优先级高于形貌分类等附加能力。

\textbf{缓解与改进。} 最高优先级是执行第\ref{sec:ground-truth}节的同批图像--机械筛分配对实验，使用校准批次拟合 $\boldsymbol{\theta},\gamma$，再在完全独立批次上验证。若低维模型在不同骨料来源间仍有稳定残差，可逐步增加形貌相关项或分材料类型校准；只有当数据量足够证明显式模型能力不足时，才考虑更高容量的学习式筛分映射。

\subsection{质量代理 $w=d^\gamma$ 的平均化假设}
\label{sec:limit-mass-proxy}

\textbf{现象与边界。} $w=d^\gamma$ 将单个尺度量映射到质量贡献，默认 $\gamma=3$ 依赖“密度近似一致、形状统计上近似相似”的平均假设。实际骨料可能在不同粒径范围内具有不同扁平程度、孔隙、岩性或厚度比例，因此相同 $d$ 的颗粒并不具有完全相同质量。

\textbf{根因。} 单张 RGB 图像没有逐颗粒厚度与密度观测。幂函数因此只能描述批次平均意义下的尺寸--质量缩放，而不是逐颗粒称重模型。

\textbf{影响。} 当形貌随尺寸系统变化时，一个全局 $\gamma$ 可能在细粒区和粗粒区产生不同方向的质量误差；即使 $D_{50}$ 经过校准，局部筛级质量比例仍可能不匹配。这也是第\ref{sec:evaluation-metrics}节要求同时报告完整筛孔曲线误差的原因。

\textbf{缓解与改进。} 第一阶段通过多个不同级配批次联合估计 $\gamma$ 并检查独立测试残差；第二阶段可在少量颗粒上采集真实质量、厚度和二维形貌，分析是否需要把 AR、面积或厚度代理引入质量模型。若使用深度/多视角获得第三维，则可从纯经验幂律进一步过渡到体积估计。

\subsection{实例分割的 merge、split 与领域迁移}
\label{sec:limit-segmentation}

\textbf{现象与边界。} 当前路线直接利用 ImageGrains 2.0 / Cellpose-SAM 预训练模型，避免现场重新训练，但预训练数据不能覆盖所有混凝土骨料表面、湿度、颜色、阴影和堆积方式。对于强阴影连接、相似背景、高反光、严重接触或异常细碎纹理，仍可能出现 merge、split、漏检和假阳性。

\textbf{根因。} 通用实例分割模型依赖训练分布中的视觉统计和边界先验。当现场域与训练域差异较大时，零样本泛化能力存在上限；此外，单张图像的模糊和遮挡也会降低任何模型可利用的边缘证据。

\textbf{影响。} split 往往制造多个偏小实例，增加细粒数量；merge 则制造一个偏大实例，并因质量权重快速放大其影响。相比普通像素 IoU，merge/split 对最终级配具有更直接的物理后果，因此应在最终视觉评估中单独统计。

\textbf{缓解与改进。} 首先用人工实例掩膜建立困难样本库，并按光照、接触、遮挡和粒径范围分层分析错误；如果错误主要集中在可通过成像改进解决的条件，优先优化采集。只有当标准化成像下仍存在稳定的领域偏差时，再使用少量本赛题标注数据进行微调或困难样本增量训练，从而避免在没有证据时盲目更换模型。

\subsection{小颗粒与分辨率下限}
\label{sec:limit-small-particles}

\textbf{现象与边界。} 当颗粒投影只覆盖很少像素时，轮廓量化误差、抗锯齿、压缩噪声和最小实例过滤都会占据相对更大比例。当前演示数据中 $<5$~mm 实例数量占比很高，但其中可能同时包含真实细颗粒和分割碎片，单靠当前规则不能完全区分。

\textbf{根因。} 图像空间分辨率决定了最小可稳定测量的边界细节；模型的 min-size 等后处理还会进一步形成算法尺度下限。像素数量不足时，即使分类为一个独立实例，轴长的相对误差也会显著增大。

\textbf{影响。} 小目标误差主要影响数量分布和 $D_{10}$ 附近，但如果实际细料质量比例较高，也会传播到完整筛分曲线。另一方面，把所有小实例一律删除又可能系统性漏掉真实 5~mm 附近颗粒。

\textbf{缓解与改进。} 通过不同相机高度/分辨率下的已知尺寸颗粒测试，建立“粒径--像素覆盖--测量误差”曲线，并据此选择现场分辨率和 min-size，而不是凭经验设置过滤阈值。若比赛只关注 5--40~mm 范围，应重点保证 5~mm 附近仍具有足够像素覆盖，并把低于有效测量下限的目标单独标记为低置信度而非简单解释为真实细料。

\subsection{二维投影形貌不能替代三维针片状判定}
\label{sec:limit-morphology}

\textbf{现象与边界。} 当前形貌模块基于 AR、圆度和 solidity 将实例分为针片状候选、圆形、棱角状和普通，但这些量全部来自俯视轮廓。一个俯视近圆的扁平颗粒可能具有很小厚度，而一个俯视细长颗粒未必满足规范意义下的片状条件。

\textbf{根因。} 真实针片状通常涉及三个正交尺寸之间的关系，单目俯视缺少厚度信息。几何阈值只能描述投影形貌，而不能恢复未观测第三维。

\textbf{影响。} 如果报告把“针片状候选”直接写成“针片状含量”，会产生定义层面的过度声明，且难以与标准量规结果比较。

\textbf{缓解与改进。} 当前报告始终保留“投影形貌/候选”表述；若竞赛需要定量评价真实针片状，应采集量规或卡尺三维尺寸作为标签，并考虑侧视、多视角或深度信息。投影规则仍可作为快速预筛，但其性能必须相对于独立物理标签报告。

\subsection{泥团与非骨料异常仍缺少材质语义}
\label{sec:limit-mud}

\textbf{现象与边界。} 当前疑似泥团规则仅使用粒径区间和 solidity。该规则可以发现部分大尺寸且明显不规则的目标，但无法区分外形近似骨料的泥团，也可能把形状天然凹凸的正常碎石误报。

\textbf{根因。} 泥团/异物的本质是材质或语义差异，而非单纯几何差异。RGB 图像中的颜色、纹理、局部表面结构可以提供更多信息，但当前规则模块尚未使用这些特征，也没有独立泥团训练/测试数据。

\textbf{影响。} 泥团漏检影响比赛异常检测任务；若把疑似泥团直接从质量级配中剔除，还可能错误删除正常粗骨料并改变 $D_p$。因此当前实现不应在没有独立确认的情况下把几何疑似标签等同于真实泥团。

\textbf{缓解与改进。} 建立正常骨料、泥团及常见异物的实例 crop 数据集，由人工材质标签监督轻量图像分类器，并将其作为实例分割后的语义分支。新分类器应独立报告 Precision/Recall，尤其关注泥团漏检和正常骨料误删两种错误；只有经过验证的异常标签才进入最终级配剔除策略。

\subsection{样本代表性与当前物理证据不足}
\label{sec:limit-evidence}

\textbf{现象与边界。} 当前仓库只提交三张具有代表性的自然堆积演示图及其推理产物，能够展示流程和行为差异，但不能覆盖不同骨料来源、含水率、光照、级配、形貌和异常组成，更没有与三张图配对的逐筛机械称量真值。

\textbf{根因。} 项目当前阶段首先完成了算法和工程闭环，而物理验证数据采集尚未同步完成。视觉任务中增加无真值图片很容易，但比赛最终评价是机械筛分语义，缺少配对物理测量就无法完成核心校准与独立验证。

\textbf{影响。} 当前最重要的不确定性不是“程序能不能运行”，而是筛分等效粒径和质量代理在真实标准筛分上的系统误差大小。因此，任何基于三张 demo 的最终准确率 claim 都是不成立的。

\textbf{缓解与改进。} 后续数据采集首先按\textbf{物理批次}而非图片数量设计：每批完成标准化拍摄与机械筛分，保存逐筛质量，再划分开发/校准/独立测试。只有在这一步完成后，才有必要继续优化 $\boldsymbol{\theta},\gamma$ 或比较更复杂模型。

\subsection{改进路线与优先级}
\label{sec:roadmap}

上述限制并不具有相同优先级。按照“对比赛主评分的影响 × 当前证据缺口 × 实施成本”，本文将后续工作排序如表\ref{tab:roadmap-priority}所示。第一优先级集中在建立物理校准闭环，因为它直接决定最终标准筛分误差；第二优先级是控制尺度和实例感知误差，使下游校准建立在可靠测量上；第三优先级才是形貌/泥团等附加语义能力。

\begin{table}[htbp]
\centering
\caption{后续改进路线及优先级。}
\label{tab:roadmap-priority}
\small
\renewcommand{\arraystretch}{1.15}
\begin{tabular}{>{\centering\arraybackslash}p{1.1cm}>{\raggedright\arraybackslash}p{3.6cm}>{\raggedright\arraybackslash}p{5.0cm}>{\raggedright\arraybackslash}p{3.9cm}}
\toprule
优先级 & 工作 & 需要新增的证据/实现 & 预期解决问题 \\
\midrule
P0 & 同批图像--机械筛分数据 & 多批逐筛称量；校准/测试严格分离 & 确定 $\boldsymbol{\theta},\gamma$，给出真实筛分误差 \\
P0 & 人工实例标注基准 & 接触/遮挡困难区域实例 mask & 量化 merge/split，分解上游误差 \\
P1 & 多点尺度与畸变校正 & 相机内参、平面标记、单应校正 & 降低场景级系统尺度偏差 \\
P1 & 困难场景适配 & 分层错误库；必要时少量微调 & 提升湿料、阴影、严重接触下的实例稳定性 \\
P2 & 泥团/异物语义分类 & 带人工材质标签的 crop 数据集 & 从几何疑似升级为可验证异常识别 \\
P2 & 三维形貌补充 & 量规/卡尺、多视角或深度 & 从投影候选升级到真实针片状评价 \\
\bottomrule
\end{tabular}
\end{table}

总体而言，本文当前方法最值得保留的并不是某个固定参数值，而是清晰的可校准结构：上游感知负责把图像转换为可检查实例，中间物理尺度把像素转换为毫米，下游低维模型把二维视觉测量映射到机械筛分语义。这样每一种新增数据都能作用于明确的误差环节，而无需把所有问题混在一个端到端网络中。下一章据此总结本文已经完成的工作、当前能够成立的结论以及在获得物理真值后应如何形成最终竞赛版本。
